% Option-A replacement fragment for the Related Work section.
% Requires citation keys from references.bib + docs/REFERENCES_2026_ADDENDUM.bib.

\section{Related Work}\label{sec:related}

\subsection{Cross-Chain Bridge Security Analysis}

Cross-chain bridge security has been studied from architectural, static-analysis, monitoring, and transaction-analysis perspectives. Prior systematizations characterize bridge architectures, recurring attack surfaces, and trust assumptions across interoperability designs~\citep{augusto2024sok,lee2023sok,duan2023attacks,notland2026sok}. These studies are particularly relevant to BridgeSentry because security properties are not universal across bridge families: authorization, settlement, refund, and finality semantics depend on the architecture and trust model of the target protocol.

Static and symbolic techniques reason about bridge code without executing an end-to-end cross-chain attack trace. SmartAxe constructs cross-chain data-flow representations and performs fine-grained static analysis for cross-chain vulnerabilities~\citep{liao2024smartaxe}. The symbolic-dataflow system of \citet{zhou2025bridgeguard} targets external-interaction vulnerabilities in bridge routers. Monitoring-oriented systems instead reason over observed execution. XScope defines cross-chain security properties over transaction/event behavior~\citep{zhang2022xscope}, HighGuard monitors cross-chain business-logic conformance~\citep{eshghie2024highguard}, and XChainWatcher identifies anomalous cross-chain behavior~\citep{augusto2025xchainwatcher}. Graph-mining systems such as the transaction-graph BridgeGuard of \citet{wu2025safeguarding} and the BridgeShield preprint~\citep{lin2025bridgeshield} model observed cross-chain transactions. These approaches address detection or conformance from source code, specifications, or observed executions; they do not by themselves establish the same experimental task as guided exploit reconstruction in a controlled execution harness.

BridgeFuzz is the closest published dynamic-testing work to BridgeSentry. It explicitly fuzzes cross-chain bridge deployments with both on-chain and off-chain relayer scope~\citep{winkler2026bridgefuzz}. Consequently, BridgeSentry does not claim novelty for cross-chain fuzzing or relayer-aware execution alone. The intended distinction is the use of an explicit source--relay--destination semantic representation, historical-incident guidance with leakage provenance, and a fail-closed validation boundary between guidance and exploit evidence. BridgeFuzz is therefore the primary dynamic baseline on the subset for which both systems can be configured faithfully.

A particularly important recent comparison is IntentFuzz~\citep{augusto2026intentfuzz}, a 2026 preprint focused on intent-based cross-chain bridges. IntentFuzz automatically recovers intent structure and deposit/fill roles from Solidity, constructs multi-step transaction sequences, checks intent-specific invariants, and uses an LLM fallback for inputs that deterministic construction cannot satisfy. Its evaluation includes labelled semantic-recovery data, planted mutants with secure controls, and forked real deployments. This work raises the empirical bar for claims about automated bridge-semantic recovery and invariant-guided cross-chain testing. BridgeSentry targets a different abstraction boundary: rather than specializing to the intent deposit/fill model, it represents source, relay, and destination dependencies through an ATG and evaluates whether historical guidance remains useful after target-specific knowledge is excluded. The final evaluation must demonstrate this distinction rather than assume it from architectural differences.

\subsection{Smart-Contract and Protocol Fuzzing}

Stateful smart-contract fuzzing has developed increasingly sophisticated mechanisms for state exploration, directed search, and exploit objectives. ItyFuzz introduced snapshot-based state exploration~\citep{shou2023ityfuzz}; SmartShot later proposed mutable snapshots for improving exploration efficiency~\citep{liang2025smartshot}; and VulSEye combines static target identification with stateful directed graybox fuzzing~\citep{liang2025vulseye}. Midas and subsequent profitable-vulnerability fuzzers optimize toward economically meaningful exploit behavior rather than coverage alone~\citep{ye2024midas,kong2025fuzzing}. Execution Property Graphs provide another representation for connecting execution behavior to security analysis~\citep{qin2025epg}. Empirical studies of smart-contract fuzzers also emphasize that repeated stochastic runs, comparable budgets, and explicit treatment of non-triggers are necessary for credible effectiveness claims~\citep{wu2024fuzzers,klees2018evaluating}.

Domain-specific guidance predates the use of language models. FuzzFactory formalized waypoint-based domain-specific feedback for fuzzing~\citep{padhye2019fuzzfactory}. In protocol fuzzing, \citet{meng2024llmguided} showed how natural-language semantics can be translated into guidance for stateful exploration. For smart contracts, LLM4Fuzz applies language-model guidance to fuzzing with user-defined invariants~\citep{shou2024llm4fuzz}. EchoFuzz further combines static analysis, LLM reasoning, vulnerable function-call sequences, and runtime feedback for smart-contract fuzzing~\citep{li2026echofuzz}. These systems mean that neither semantic waypoints nor ``LLM + fuzzing'' is sufficient as a novelty claim for BridgeSentry. The relevant question is whether bridge-specific cross-domain semantics and leakage-controlled historical guidance improve valid exploit traces under a fixed harness and oracle.

\subsection{LLM- and Retrieval-Assisted Security Analysis}

LLMs have been used for smart-contract vulnerability detection, property generation, and verification. GPTScan combines GPT-based semantic reasoning with program analysis~\citep{sun2024gptscan}, while multi-agent approaches coordinate multiple LLM roles for vulnerability analysis~\citep{wei2025advanced}. SmartInv learns smart-contract invariants from multiple modalities~\citep{wang2024smartinv}, and INVCON+ combines invariant inference with validation mechanisms~\citep{liu2025invcon}. These works motivate treating an LLM-generated invariant as a candidate rather than as ground truth.

Retrieval-augmented security reasoning is likewise established prior art. PropertyGPT retrieves relevant properties to support LLM-driven formal verification~\citep{liu2025propertygpt}. ParaVul combines an LLM vulnerability detector with hybrid sparse/dense retrieval and gated verification~\citep{huang2026paravul}. BridgeSentry therefore does not claim that retrieval augmentation itself is novel. Its evaluation instead treats the historical exploit corpus as a potential source of target leakage: every generated scenario retains retrieved incident identifiers, and effectiveness is evaluated separately under full-knowledge, leave-one-incident-out, leave-one-bridge-family-out, and vulnerability-class/temporal holdout regimes. This separation is necessary because successful reconstruction with access to a target incident does not establish generalization.

LLM pretraining introduces a second contamination channel that retrieval exclusion cannot remove. Public exploit postmortems, source code, and attack reports may be present in model pretraining data. For this reason, model/provider identifiers and dates must be retained in the experiment manifest, and parser-only or non-LLM ablations are used where feasible to bound the contribution of parametric knowledge.

\subsection{Atomic Transfer Graphs and BridgeSentry's Position}

Atomic Transfer Graphs (ATGs) were introduced as a secure-by-design abstraction for heterogeneous blockchain protocols~\citep{dubler2025atg}. BridgeSentry adapts the representation for a different purpose: extracting a test-oriented intermediate representation from existing bridge artifacts and associating it with protocol-specific candidate security properties. The adaptation is useful only if the extracted representation is independently evaluated; self-generated ATG coverage is therefore treated as an internal progress measure rather than as evidence of semantic correctness.

Taken together, the literature leaves BridgeSentry with a narrower claim than an earlier ``LLM-guided cross-chain fuzzer'' framing. Cross-chain dynamic testing is established by BridgeFuzz and IntentFuzz, LLM-guided fuzzing is established by prior protocol and smart-contract work, and retrieval-assisted security reasoning is also established. BridgeSentry's target contribution is the integration of (i) an explicit source--relay--destination semantic intermediate representation, (ii) provenance-aware and leakage-controlled historical guidance, and (iii) execution-grounded, fail-closed exploit validation. The experiments are designed to test each of these elements independently rather than infer their value from a target-aware predicate-trigger result.
